\section{What does the BERT router attend to?}
\label{sec:interpretability}

To understand what makes $\mathcal{R}_{\mathrm{BERT}}$ prefer Lean, we analyse the test steps it routes to Lean with the default threshold, i.e.\ the 12 steps with $P(\mathrm{Lean})\geq0.5$ (all between 50.0\% and 54.0\%, so near-boundary rather than confident decisions). For each step we attribute the Lean logit to the input tokens with Integrated Gradients \cite{sundararajan2017axiomatic} and then test the most attributed words directly, by deleting or masking one occurrence and re-scoring the step (details in Appendix~\ref{app:ig}). Figure~\ref{fig:lean_triggers} shows three of the steps and the effect of the interventions.

\begin{figure*}[t]
    \centering
    \IfFileExists{figures/lean_triggers_heatmaps.pdf}{\includegraphics[width=\textwidth]{lean_triggers_heatmaps.pdf}}{\fbox{\texttt{figures/lean\_triggers\_heatmaps.pdf} missing}}\\[2pt]
    \IfFileExists{figures/lean_triggers_interventions.pdf}{\includegraphics[width=\textwidth]{lean_triggers_interventions.pdf}}{\fbox{\texttt{figures/lean\_triggers\_interventions.pdf} missing}}
    \caption{Top: word-level Integrated Gradients attribution of the Lean logit for three of the 12 test steps routed to Lean (blue supports, orange opposes Lean; scores are $L_2$-normalised within each step; E1 has routing label Lean, E2 and E3 Python). Bottom: reduction of $P(\mathrm{Lean})$, in percentage points, when the most attributed word of each step (left) or a low-attribution comparison word (right) is deleted (circles) or masked (squares).}
    \label{fig:lean_triggers}
\end{figure*}

The recurring cue is qualifying and negative language rather than mathematical content: \emph{however} contributes to the Lean decision in 6 of the 12 steps, and \emph{doesn't}, \emph{isn't}, \emph{not}, \emph{can't} and \emph{unclear} in several others, whereas the comparison words move it by less than one percentage point. Removing a single such word is enough to move some steps back across the default threshold (deleting \emph{doesn't} in E2 lowers $P(\mathrm{Lean})$ from 51.9\% to 47.1\%, masking it to 43.0\%), and removing the five most attributed words jointly lowers $P(\mathrm{Lean})$ in most steps, but never below 34\%, far above the validation-selected threshold of 3.4\%. The router has thus learned that steps expressing doubt, contrast or a negative claim are the ones on which the Python verdict tends to disagree with the annotation, a cue that is indicative of the \emph{kind} of step rather than of whether a given verifier will succeed on it, and that is consistent with the small downstream gain of the learned router over always choosing Python (Table~\ref{tab:test_split_results}). The same picture holds over the whole test split: attributing the routing margin of all 1,363 test steps (Appendix~\ref{app:ig}), negation words have by far the largest mean attribution towards Lean (\emph{not}, \emph{however}, \emph{doesn't}), whereas the words that most push towards Python are the connectives and verbs of routine calculation (\emph{therefore}, \emph{thus}, \emph{let's}, \emph{calculate}).
